% Scaffold copied from thesis/planned-toc.md.
% This file intentionally contains headings and local TODO/context comments only.
% Old thesis prose lives in thesis/legacy/ as source material.

\section{Algorithm For First-Order Perturbations}\label{sec:first-order-perturbations}
%
% Purpose:
% explain how local perturbations of a polytope are evaluated for gradient-like
% search and for the HKO2024 local-maximum computation.
%
\subsection{Definition}\label{subsec:first-order-perturbations-definition}
%
% Define the first-order objects in the notation used by the rest of the thesis.
% <!-- OUTLINE GAP: Name the objects: row-coordinate perturbation h, active
%      orbit/word set, branch action, beta derivative or subgradient, volume
%      derivative, and derivative of sys. -->
%
% Writing note:
% introduce the generic case first because it is readable. Handle edge cases only
% where the HKO computation or the thesis claims need them.
% <!-- OUTLINE GAP: Decide the exact generic hypotheses to list here:
%      positive dwell times, full-rank constraints, negative-definite reduced
%      Hessian, unique active minimizer or finite active set, and fixed face
%      combinatorics for volume. -->
% <!-- WRITER-SESSION QUESTION: Decide while writing the chapter which
%      non-generic cases matter for HKO2024. Do not attempt to cover every
%      a-priori possible degeneracy unless it appears in, or is needed for, the
%      HKO argument. -->
%
\subsection{Correctness}\label{subsec:first-order-perturbations-correctness}
%
% Prove the generic case cleanly. Add non-generic cases afterward as needed for
% HKO2024 or for honest caveats.
% <!-- OUTLINE GAP: Split into theorem statements: differentiability of one
%      branch, derivative formula for capacity/action, derivative formula for
%      volume, derivative/subgradient statement for sys, and separate caveat for
%      ties or semidefinite cases. -->
%
\subsection{Empirical tests}\label{subsec:first-order-perturbations-empirical-tests}
%
% State the numerical checks that the implementation behaves as expected on the
% examples used in the thesis.
% <!-- OUTLINE GAP: Name the checks: finite-difference agreement, active-set
%      stability, comparison to ascent behavior, and HKO-specific row/rank
%      bookkeeping. -->
%
